% --- コンパイル環境 ---
% コンパイラ: LuaLaTeX
% 文字コード: UTF-8

\documentclass[11pt,dvipsnames]{beamer}
\usepackage{luatexja} % LuaLaTeX用日本語パッケージ
\usepackage{amsmath, amssymb, amsthm}

% テーマとカラー設定
\usetheme{Madrid}
\usecolortheme{seahorse}

% 定理環境の設定
\setbeamertemplate{theorems}[numbered]
\renewcommand{\proofname}{証明}
\newtheorem{thm}{定理}
\newtheorem{lem}[thm]{補題}
\newtheorem{dfn}[thm]{定義}
\newtheorem{ex}[thm]{例}

% スライド下部のナビゲーションシンボルを非表示
\setbeamertemplate{navigation symbols}{}

% 番号付きリストの設定
\setbeamertemplate{enumerate items}[default]

% タイトル情報
\title[Galois拡大の4条件]{有限次Galois拡大における\\4つの同値条件の完全証明}
\subtitle{自己完結版}
\author{}
\date{\today}

\begin{document}

% --- タイトルスライド ---
\begin{frame}
    \titlepage
\end{frame}

% --- 目次スライド ---
\begin{frame}{目次}
    \tableofcontents
\end{frame}

% ==========================================
\section{1. 基本概念の定義と具体例}
% ==========================================

\begin{frame}[allowframebreaks]{1. 基本概念の定義と具体例}
    本稿を通じて用いる基本的な代数的概念を厳密に定義する。以下、断りのない限り $L/K$ は体の拡大を表し、$[L:K]$ は $K$ 上のベクトル空間としての $L$ の次元（拡大次数）を表す。本稿では有限次拡大、すなわち $[L:K] < \infty$ の場合を扱う。また、$\mathrm{Aut}(L/K)$ は $K$ の各元を固定する $L$ の体自己同型全体のなす群を表す。

    \begin{dfn}[基本概念の定義]
        \begin{itemize}
            \item \textbf{代数拡大 (algebraic extension)}: 拡大 $L/K$ の任意の元 $\alpha \in L$ に対し、$\alpha$ を根に持つような $K$ 上の非零多項式 $f(x) \in K[x]$ が存在するとき、$L/K$ は代数拡大であるという。
            \item \textbf{最小多項式 (minimal polynomial)}: 代数的な元 $\alpha \in L$ に対し、$\alpha$ を根に持つ $K[x]$ のモニックな既約多項式を $\alpha$ の $K$ 上の最小多項式と呼ぶ。
            \item \textbf{分離的 (separable)}: 既約多項式 $f(x) \in K[x]$ がその代数閉包 $\overline{K}$ において重根を持たないとき、$f(x)$ は分離多項式であるという。元 $\alpha \in L$ の最小多項式が分離多項式であるとき、$\alpha$ は $K$ 上分離的であるという。$L$ のすべての元が $K$ 上分離的であるとき、拡大 $L/K$ は分離拡大であるという。
            \item \textbf{正規拡大 (normal extension)}: 既約多項式 $f(x) \in K[x]$ が $L$ に少なくとも1つの根を持つならば、$L[x]$ において一次式の積に完全に分解するとき、拡大 $L/K$ は正規拡大であるという。
            \item \textbf{最小分解体 (splitting field)}: 多項式 $f(x) \in K[x]$ に対し、ある拡大体 $L/K$ が $f(x)$ のすべての根 $\alpha_1, \dots, \alpha_n$ を含み、かつ $L = K(\alpha_1, \dots, \alpha_n)$ と表されるとき、$L$ を $f(x)$ の $K$ 上の最小分解体と呼ぶ。
            \item \textbf{不変体 (fixed field)}: 自己同型群の有限部分群 $G \subset \mathrm{Aut}(L)$ に対し、$L^G = \{ x \in L \mid \forall \sigma \in G, \sigma(x) = x \}$ と定義される体を $G$ の不変体と呼ぶ。
        \end{itemize}
    \end{dfn}
\end{frame}

\begin{frame}[allowframebreaks]{正規性と分離性の具体例}
    \begin{ex}[正規性と分離性の具体例]
        \begin{itemize}
            \item $L = \mathbb{Q}(\sqrt[3]{2})$ と $K = \mathbb{Q}$ を考える。$\sqrt[3]{2}$ の $K$ 上の最小多項式は $x^3 - 2$ である。この多項式は $\mathbb{Q}$ 上既約であり、$\overline{\mathbb{Q}}$ における根は相異なる3つの複素数となるため分離的であるが、虚数根が $L$ に含まれないため、$L/K$ は正規拡大ではない。
            \item 標数 $p > 0$ の素体 $\mathbb{F}_p$ 上の一変数有理関数体 $K = \mathbb{F}_p(t^p)$ と、その拡大体 $L = \mathbb{F}_p(t)$ を考える。$t \in L$ の $K$ 上の最小多項式は $x^p - t^p$ であるが、これは $(x-t)^p$ と因数分解されるため、$t$ を $p$ 重根として持つ。したがって、この拡大は分離拡大ではない。
        \end{itemize}
    \end{ex}
\end{frame}

% ==========================================
\section{2. 前提となる定理・補題の証明}
% ==========================================

\begin{frame}[allowframebreaks]{2. 前提となる定理・補題の証明}
    同値性の証明に先立ち、必要となるすべての基本定理を完全に証明し、自己完結した議論を行う。

    \begin{lem}[Dedekindの写像の線形独立性]
        $L, M$ を任意の体とする。$L$ から $M$ への相異なる体準同型 $\sigma_1, \dots, \sigma_n$ は、$M$ 上線形独立である。すなわち、$\lambda_1, \dots, \lambda_n \in M$ に対して、すべての $x \in L$ で $\sum_{i=1}^n \lambda_i \sigma_i(x) = 0$ が成り立つならば、$\lambda_1 = \dots = \lambda_n = 0$ である。
    \end{lem}
    
    \begin{proof}
        非零の係数を持つ線形従属な関係式が存在すると仮定し、項数 $k$ が最小のものを $\sum_{i=1}^k \lambda_i \sigma_i(x) = 0$ とする。$\sigma_1 \neq \sigma_k$ より、ある $\alpha \in L$ が存在して $\sigma_1(\alpha) \neq \sigma_k(\alpha)$ となる。任意数 $x \in L$ に対し、上の式に $\alpha x$ を代入すると、
        \[ \sum_{i=1}^k \lambda_i \sigma_i(\alpha) \sigma_i(x) = 0 \]
        一方で、元の関係式に $\sigma_k(\alpha)$ を乗じると、
        \[ \sum_{i=1}^k \lambda_i \sigma_k(\alpha) \sigma_i(x) = 0 \]
        これら2式の差をとると、第 $k$ 項が相殺され、
        \[ \sum_{i=1}^{k-1} \lambda_i (\sigma_i(\alpha) - \sigma_k(\alpha)) \sigma_i(x) = 0 \]
        $i=1$ のとき $\lambda_1(\sigma_1(\alpha) - \sigma_k(\alpha)) \neq 0$ であるため、これは項数が $k$ 未満の自明でない線形従属関係式となり、最小性に矛盾する。
    \end{proof}
\end{frame}

\begin{frame}[allowframebreaks]{Artinの定理}
    \begin{thm}[Artinの定理]
        体 $L$ と、$\mathrm{Aut}(L)$ の有限部分群 $G = \{ \sigma_1 = 1, \sigma_2, \dots, \sigma_n \}$ を考える。$K = L^G$ とおくと、拡大次数について $[L:K] = n = |G|$ が成立する。
    \end{thm}
    
    \begin{proof}
        \textbf{(Step 1) $[L:K] \ge n$ の証明}\\
        $[L:K] = m < n$ と仮定する。$K$ 上の $L$ の基底を $\omega_1, \dots, \omega_m$ とする。次のような方程式系を考える：
        \[ \sum_{j=1}^n \sigma_j(\omega_i) x_j = 0 \quad (i = 1, \dots, m) \]
        式の数 $m$ が未知数の数 $n$ より少ないため、非自明な解 $(x_1, \dots, x_n)$ を持つ。任意の $\alpha = \sum_{i=1}^m a_i \omega_i \in L$ に対し、$a_i \in K$ を乗じて $i$ について和をとると、
        \[ \sum_{j=1}^n \sigma_j(\alpha) x_j = 0 \]
        これは体自己同型が $L$ 上線形従属であることを意味し、補題 1 に矛盾する。したがって $[L:K] \ge n$。

        \textbf{(Step 2) $[L:K] \le n$ の証明}\\
        $[L:K] > n$ と仮定する。$K$ 上線形独立な $\alpha_1, \dots, \alpha_{n+1} \in L$ が存在する。方程式系 $\sum_{i=1}^{n+1} \sigma_j(\alpha_i) y_i = 0$ は非自明な解を持つ。非零成分の個数が最小となる解を選び、$y_r = 1$ としてよい。
        方程式系に $\sigma_k \in G$ を作用させ、元の式との差をとると、さらに項数の少ない関係式が得られ、解の最小性に矛盾する。したがって $[L:K] \le n$。以上より $[L:K] = n$ である。
    \end{proof}
\end{frame}

\begin{frame}[allowframebreaks]{埋め込みの数と根の数}
    \begin{lem}[単一拡大における埋め込みの数と根の数]
        $K(\alpha)/K$ を単一拡大とし、$\alpha$ の $K$ 上の最小多項式を $p(x)$ とする。$K$ を固定する埋め込み $\sigma: K(\alpha) \to \overline{K}$ の総数は、$p(x)$ の $\overline{K}$ における相異なる根の数に等しい。また、この総数が $[K(\alpha):K]$ と一致することと、$\alpha$ が $K$ 上分離的であることは同値である。
    \end{lem}
    
    \begin{proof}
        $K(\alpha)$ の任意の元は $g(\alpha)$ と表せる。$\sigma$ の効果は $\sigma(\alpha)$ の行き先だけで決定される。$0 = \sigma(p(\alpha)) = p(\sigma(\alpha))$ より、$\sigma(\alpha)$ は $p(x)$ の根でなければならない。逆に $p(x)$ の任意の根 $\beta$ に対し、$\alpha \mapsto \beta$ とする準同型が定まる。埋め込みの総数は相異なる根の数に等しく、これが次数と一致するのは重根を持たない（分離的）ことの定義そのものである。
    \end{proof}
\end{frame}

\begin{frame}[allowframebreaks]{分離次数の性質}
    \begin{thm}[分離次数の性質と分離元生成の拡大]
        有限次拡大 $L/K$ に対し、$\overline{K}$ への $K$ 準同型の総数を分離次数 $[L:K]_s$ と定義する。
        \begin{enumerate}
            \item \textbf{乗法性}: 中間体 $M$ に対し $[L:K]_s = [L:M]_s [M:K]_s$
            \item \textbf{次数不等式}: $[L:K]_s \le [L:K]$ であり、等号成立は分離拡大と同値。
            \item \textbf{分離元による生成}: $\alpha_1, \dots, \alpha_r$ がすべて分離的であれば、$K(\alpha_1, \dots, \alpha_r)/K$ も分離拡大。
        \end{enumerate}
    \end{thm}
    
    \begin{proof}
        1. 埋め込みの拡張と合成から明らか。\\
        2. 単一拡大の塔に分解し、各段階で補題 3 を適用。等号成立は各段階の生成元が分離的であることと同値。\\
        3. $\alpha_2$ の $K$ 上の最小多項式 $f(x)$ は重根を持たない。$M=K(\alpha_1)$ 上の最小多項式 $g(x)$ は $f(x)$ を割り切るため、やはり重根を持たない。乗法性と 2 より $K(\alpha_1, \alpha_2)/K$ は分離拡大。帰納法により成立。
    \end{proof}
\end{frame}

\begin{frame}[allowframebreaks]{最小分解体の正規性}
    \begin{thm}[最小分解体の正規性]
        $L$ が多項式 $f(x) \in K[x]$ の $K$ 上の最小分解体であるならば、$L/K$ は正規拡大である。
    \end{thm}
    
    \begin{proof}
        $f(x)$ の $L$ における根全体を $\alpha_1, \dots, \alpha_m$ とすると $L = K(\alpha_1, \dots, \alpha_m)$。任意の $K$ 埋め込み $\sigma$ は根を置換するため $\sigma(L) = L$。$g(x) \in K[x]$ を $L$ に根 $\beta$ を持つ既約多項式とする。$\overline{K}$ の他の根 $\gamma$ への埋め込みを $L$ に拡張すると $\sigma(L) = L$ より $\gamma \in L$ となる。よって正規拡大。
    \end{proof}
\end{frame}

% ==========================================
\section{3. 4つの条件の同値性の直接証明}
% ==========================================

\begin{frame}{4つの条件}
    以下の4つの条件が互いに同値であることを、12通りの含意について直接証明する。
    
    \begin{block}{同値な4条件}
        \begin{enumerate}
            \item[(1)] $L/K$ は分離的正規拡大である。
            \item[(2)] $L^G = K$ である。
            \item[(3)] $[L:K] = |G|$ である。
            \item[(4)] $L$ は重根を持たないある既約多項式 $f(x) \in K[x]$ の $K$ 上での最小分解体である。
        \end{enumerate}
    \end{block}
\end{frame}

% --- (1) からの導出 ---

\begin{frame}[allowframebreaks]{(1) $\Rightarrow$ (2) の証明}
    \textbf{【分離次数を使わない証明】}\\
    $L^G \subset K$ を示す。任意の $\alpha \in L \setminus K$ をとり、最小多項式を $p(x)$ とする。分離性より重根を持たず、正規性より $L$ で分解する。$\alpha \notin K$ より $\deg p \ge 2$ なので、異なる根 $\beta \in L$ が存在する。$\alpha \mapsto \beta$ となる $L$ の自己同型 $\sigma \in G$ が存在し $\sigma(\alpha) \neq \alpha$ となるため、$\alpha \notin L^G$。よって $L^G = K$。

    \vspace{1em}
    \textbf{【分離次数を使う証明】}\\
    分離性より $[L:K]_s = [L:K]$。正規性より $[L:K]_s = |G|$。ゆえに $[L:K] = |G|$。Artinの定理より $[L:L^G] = |G|$。次数の連鎖律 $[L:K] = [L:L^G][L^G:K]$ より $[L^G:K] = 1$。よって $L^G = K$。
\end{frame}

\begin{frame}[allowframebreaks]{(1) $\Rightarrow$ (3) の証明}
    \textbf{【分離次数を使わない証明】}\\
    既に (1) $\Rightarrow$ (2) により $L^G = K$ が示されている。Artinの定理より $[L:L^G] = |G|$ であるから、直ちに $[L:K] = |G|$ を得る。

    \vspace{1em}
    \textbf{【分離次数を使う証明】}\\
    分離性より $[L:K]_s = [L:K]$、正規性より $[L:K]_s = |G|$。これらを直接結ぶことで、$[L:K] = |G|$ が直ちに得られる。
\end{frame}

\begin{frame}[allowframebreaks]{(1) $\Rightarrow$ (4) の証明}
    \textbf{【分離次数を使わない証明】}\\
    $L = K(\alpha_1, \dots, \alpha_m)$ と表す。各 $\alpha_i$ の最小多項式 $p_i(x)$ は、分離性より重根を持たず、正規性より $L$ で分解する。$p_1(x) \cdots p_m(x)$ の無平方部分を $f(x)$ とすれば、$L$ は $f(x)$ の最小分解体となる。

    \vspace{1em}
    \textbf{【分離次数を使う証明】}\\
    分離性より（定理2.4-2から）各元の最小多項式は重根を持たない。正規性により $L$ で分解する。同様に無平方部分 $f(x)$ をとることで構成できる。
\end{frame}

% --- (2) からの導出 ---

\begin{frame}[allowframebreaks]{(2) $\Rightarrow$ (1) の証明}
    \textbf{【分離次数を使わない証明】}\\
    任意の $\alpha \in L$ の軌道を $S = \{ \sigma(\alpha) \mid \sigma \in G \} = \{ \alpha_1, \dots, \alpha_k \}$ とし、$g(x) = \prod_{i=1}^k (x - \alpha_i)$ を考える。係数は $G$ 不変であり、(2) より $L^G = K$ なので $g(x) \in K[x]$。構成から $g(x)$ は相異なる根のみを持ち完全に分解する。$\alpha$ の最小多項式は $g(x)$ を割り切るため、分離的かつ正規である。

    \vspace{1em}
    \textbf{【分離次数を使う証明】}\\
    Artinの定理 $[L:L^G] = |G|$ に $L^G = K$ を代入し $[L:K] = |G|$。一般に $|G| \le [L:K]_s \le [L:K]$ であるから、すべての不等号が等号になる。$[L:K]_s = [L:K]$ より分離的、$|G| = [L:K]_s$ より正規である。
\end{frame}

\begin{frame}[allowframebreaks]{(2) $\Rightarrow$ (3) の証明}
    \textbf{【分離次数を使わない証明 / 使う証明 (共通)】}\\
    Artinの定理により $[L:L^G] = |G|$。ここに仮定 (2) である $L^G = K$ を直接代入することにより、直ちに $[L:K] = |G|$ が導かれる。
\end{frame}

\begin{frame}[allowframebreaks]{(2) $\Rightarrow$ (4) の証明}
    \textbf{【分離次数を使わない証明】}\\
    (2) $\Rightarrow$ (1) の証明により、$L/K$ が分離的かつ正規であることが示される。その後は (1) $\Rightarrow$ (4) と全く同じ手順で無平方部分 $f(x)$ を構成する。

    \vspace{1em}
    \textbf{【分離次数を使う証明】}\\
    (2) $\Rightarrow$ (1) の分離次数を使う証明により分離的正規性を得て、同様に $f(x)$ を構成する。
\end{frame}

% --- (3) からの導出 ---

\begin{frame}[allowframebreaks]{(3) $\Rightarrow$ (1) の証明}
    \textbf{【分離次数を使わない証明】}\\
    単一拡大の塔 $K_i = K_{i-1}(\beta_i)$ を考える。各段階の埋め込みの拡張数は最大で $[K_i:K_{i-1}]$ であり、重根を持たず根が $L$ に含まれるときに最大となる。$|G| \le \prod [K_i:K_{i-1}] = [L:K]$ であり、仮定 (3) より等号が成立。すべての段階で拡張可能数は最大かつ像は $L$ に含まれるため、すべての元の最小多項式は重根を持たず $L$ で分解する。

    \vspace{1em}
    \textbf{【分離次数を使う証明】}\\
    $|G| \le [L:K]_s \le [L:K]$ において仮定より両端が一致するため、$[L:K]_s = [L:K]$（分離的）および $|G| = [L:K]_s$（正規的）が同時に成立する。
\end{frame}

\begin{frame}[allowframebreaks]{(3) $\Rightarrow$ (2) の証明}
    \textbf{【分離次数を使わない証明 / 使う証明 (共通)】}\\
    次数の連鎖律 $[L:K] = [L:L^G][L^G:K]$ に、Artinの定理 $[L:L^G] = |G|$ と仮定 (3) の $[L:K] = |G|$ を代入すると、$|G| = |G| [L^G:K]$。よって $[L^G:K] = 1$ となり $L^G = K$ を得る。
\end{frame}

\begin{frame}[allowframebreaks]{(3) $\Rightarrow$ (4) の証明}
    \textbf{【分離次数を使わない証明 / 使う証明 (共通)】}\\
    (3) $\Rightarrow$ (1) により、分離的かつ正規であることが示される。あとは (1) $\Rightarrow$ (4) と同様に有限個の生成元の最小多項式から $f(x)$ を構成すればよい。
\end{frame}

% --- (4) からの導出 ---

\begin{frame}[allowframebreaks]{(4) $\Rightarrow$ (1) の証明}
    \textbf{【分離次数を使わない証明】}\\
    定理2.5より正規拡大である。分離性について、$f(x)$ の根の塔を考えると、各段階の最小多項式は $f(x)$ を割り切るため重根を持たず、埋め込み総数は $[L:K]$ となる。正規性により $|G| = [L:K]$。Artinの定理より $L^G = K$ を得る。(2) $\Rightarrow$ (1) と同様に軌道から $g(x)$ を構成することで、任意の元が分離的であることが示される。

    \vspace{1em}
    \textbf{【分離次数を使う証明】}\\
    定理2.5より正規拡大。$f(x)$ は重根を持たないため、その根は $K$ 上分離的。定理2.4-3（分離的生成元による拡大は分離的）を適用することで、$L$ は分離拡大となる。
\end{frame}

\begin{frame}[allowframebreaks]{(4) $\Rightarrow$ (2) の証明}
    \textbf{【分離次数を使わない証明】}\\
    (4) $\Rightarrow$ (1) の途中で示した通り、埋め込みの数え上げから $|G| = [L:K]$。Artinの定理 $[L:L^G] = |G|$ より $L^G = K$ を得る。

    \vspace{1em}
    \textbf{【分離次数を使う証明】}\\
    (4) $\Rightarrow$ (1) により分離的かつ正規。分離性より $[L:K]_s = [L:K]$、正規性より $[L:K]_s = |G|$。よって $[L:K] = |G|$。Artinの定理より $L^G = K$。
\end{frame}

\begin{frame}[allowframebreaks]{(4) $\Rightarrow$ (3) の証明}
    \textbf{【分離次数を使わない証明】}\\
    根を順に添加する塔において、各段階の最小多項式が重根を持たず、拡張可能数が正確に $[K_i:K_{i-1}]$ となる。これを掛け合わせることで埋め込み数が $[L:K]$ となり、正規性によりこれらがすべて自己同型となるため $|G| = [L:K]$ を得る。

    \vspace{1em}
    \textbf{【分離次数を使う証明】}\\
    (4) $\Rightarrow$ (1) により分離的正規。分離次数の性質から $[L:K] = [L:K]_s$ かつ $[L:K]_s = |G|$。これらを直結させて $[L:K] = |G|$ が得られる。
\end{frame}

% ==========================================
\section{参考文献}
% ==========================================

\begin{frame}{引用文献}
    \begin{itemize}
        \item Artin, E. (1944). \textit{Galois Theory}. Notre Dame Mathematical Lectures, no. 2. University of Notre Dame Press.
        \item Milne, J. S. (2022). \textit{Fields and Galois Theory}. Course Notes.
        \item Lang, S. (2002). \textit{Algebra} (Revised 3rd ed.). Graduate Texts in Mathematics, 211. Springer-Verlag.
    \end{itemize}
\end{frame}

\end{document}